\section{Comparing the principal insertion with its Gaussian reference}
\label{sec:principal-gaussian-comparison}
The principal cluster is not itself a small remainder. Its leading
part is a local Gaussian reference observable. We compare these two
objects on one probability space, keeping the insertion fixed while
differentiating the surrounding densities. This is necessary even
when the insertion is only bounded and measurable.

\subsection{The common reference integral}
Fix a forest, its tags, a regular insertion of amplitude $\tau\ge0$,
and a healthy base configuration. Let $y\in\mathbb R^j$, $j\le n$,
parameterize an ordered local variation about that configuration,
with base point $0$. The read coordinates form a finite set of size
$N_0$, including every colour of each read bond. Write
$C_6(y),C_9(y)$ for the actual-sector and Gaussian-reference
covariances on these coordinates, and $w_6(y,z),w_9(y,z)$ for
the two compensated coordinate maps. At the base point
\[
                    C_0=C_6(0),\qquad w_6(0,z)=z.
\]
Both covariances are positive in a neighborhood of zero.
Let $J(y,z)$ be the common compensated Jacobian and put
\[
 R_a(y,z)=J(y,z)\frac{p_{C_a(y)}(w_a(y,z))}{p_{C_0}(z)},
                                     \qquad a=6,9,
\]
where $p_C$ is the Gaussian density in
\eqref{eq:gaussian-local-density}. The reference probability
measure is always $\gamma_{C_0}$.

Let $h(z)$ be the original masked insertion evaluated at the fixed
base configuration and its original Gaussian coordinates. It is
measurable and $|h(z)|\le\tau$. The exact two integrands are
\begin{equation}\label{eq:principal-two-integrands}
 R_6(y,z)\chi_6(y,z)h(z),\qquad
 R_9(y,z)\chi_9(y,z)h(z).
\end{equation}
Here $\chi_6$ is the sector cutoff and $\chi_9$ the reference
cutoff, both smooth with values in $[0,1]$. The first identity
follows from the compensated change of variables of
Section~\ref{sec:gaussian-differentiation}: on the support of its
cutoff the new configuration agrees with the fixed one on the
insertion's read support. Locality therefore identifies their
insertion values. For the second identity, the reference cutoff
implies that the first cutoff is one and that the fixed configuration
is healthy on the insertion support. The mask can then be removed,
and the same local agreement identifies the unmasked reference
insertion with $h(z)$. Where a cutoff vanishes both sides of its
identity are zero. Thus no derivative of $h$ is taken in either
formula, and no insertion-dependent change of Gaussian measure
has been made.

We state the quantitative finite-dimensional data used below
explicitly. They are uniform over all the forests, tags and local
variations under consideration. There are $N,Q,S\ge1$,
$\delta\ge0$, and a fixed $M>0$ such that:
\begin{enumerate}
\item $N_0\le N$ and $C_0\le MI$;
\item the entries of $C_6,C_9$ have jets through order $n$
bounded by $Q$ at zero, their inverses at zero have entries
bounded by $Q$, and the jets of $C_9-C_6$ are bounded by $\delta$;
\item for $\|z\|_\infty\le S$, the entries of $w_6,w_9$ and
the scalar $J$ have jets bounded by $Q$, and those of
$w_9-w_6$ are bounded by $\delta$;
\item on a common neighborhood of zero both integrands in
\eqref{eq:principal-two-integrands} vanish for $\|z\|_\infty>S$.
The cutoff jets on this ball are bounded by a fixed polynomial
in $N,Q,S$. For every $z$ with $\|z\|_\infty<v$, both cutoffs
are identically one on some neighborhood of zero.
\end{enumerate}
Here a jet bound includes the value and the operator norms of
derivatives of all orders at most $n$. The last neighborhood may
depend on $z$; the common compact support neighborhood does not.
The covariances, coordinate maps and Jacobian are smooth on the
compact coordinate ball and a neighborhood of the base parameter.
These hypotheses are the specific comparison data; they are not
consequences of an abstract equality between the two gas profiles.

\subsection{A quantitative density difference}
There is a polynomial $P_n$ with positive coefficients, depending
only on $n$, such that, after increasing $P_n$ if necessary and
provided
\begin{equation}\label{eq:principal-delta-small}
                      \delta P_n(N,Q,S)\le1,
\end{equation}
the following bounds hold for $\|z\|_\infty\le S$:
\begin{equation}\label{eq:principal-density-jets}
 \|R_6(\cdot,z)\|_{J^n_0}\le P_n(N,Q,S),\qquad
 \|R_9(\cdot,z)-R_6(\cdot,z)\|_{J^n_0}
                                      \le\delta P_n(N,Q,S).
\end{equation}
We prove the polynomial form of these estimates rather than leave
the dependence on the number of read coordinates implicit.

For matrices use the maximum entry norm. Multiplication costs at
most $N$ per contracted index. Differentiating $CC^{-1}=I$
gives, for a nonempty labelled derivative set $I$,
\[
 D_I(C^{-1})=-C^{-1}
        \sum_{\varnothing\ne A\subseteq I}(D_AC)
                                      D_{I\setminus A}(C^{-1}).
\]
Induction in $|I|\le n$ bounds all inverse jets by a polynomial
in $N,Q$. The exact identity
\[
 C_9^{-1}-C_6^{-1}
                      =C_9^{-1}(C_6-C_9)C_6^{-1}
\]
and the Leibniz rule give a factor $\delta$ in every derivative
of the inverse difference. The remaining factors have the
polynomial bounds just obtained.

For the logarithmic determinant, its positive-order derivatives
follow by repeatedly differentiating
$D\log\det C=\operatorname{tr}(C^{-1}DC)$.
Telescoping each product shows that the difference of these
derivatives for $C_9$ and $C_6$ is $\delta$ times a polynomial.
For the zeroth order join $C_6(0)$ to $C_9(0)$ by the straight
positive segment. Since the endpoint inverse operator norms are
at most $NQ$, both endpoints are at least $(NQ)^{-1}I$;
the whole segment has the same lower bound. Hence
\[
 |\log\det C_9(0)-\log\det C_6(0)|
 \le\int_0^1N\|C_s^{-1}\|\|C_9(0)-C_6(0)\|\,ds
 \le N^3Q\delta.
\]
The two quadratic forms in the Gaussian exponent differ by
\[
\begin{split}
 w_9^TC_9^{-1}w_9-w_6^TC_6^{-1}w_6
  ={}&(w_9-w_6)^TC_9^{-1}w_9\\
     &+w_6^TC_9^{-1}(w_9-w_6)
       +w_6^T(C_9^{-1}-C_6^{-1})w_6.
\end{split}
\]
Every derivative therefore has the same $\delta$ factor.

Let $\ell_a=\log[p_{C_a(y)}(w_a(y,z))/p_{C_0}(z)]$.
The preceding calculations bound all jets of $\ell_9-\ell_6$
by $\delta$ times a polynomial. All positive-order jets of
$\ell_6$ have polynomial bounds, and $\ell_6(0,z)=0$.
The labelled-partition formula
\[
 D_Ie^\ell=e^\ell\sum_{\pi\in\operatorname{Part}(I)}
                                  \prod_{A\in\pi}D_A\ell
\]
thus gives polynomial jets of $e^{\ell_6}$. For
$e^{\ell_9-\ell_6}-1$, the value is bounded by
$e|\ell_9-\ell_6|$ under \eqref{eq:principal-delta-small}.
Each positive-order partition has at least one derivative factor,
so all its terms retain a factor $\delta$. The other factors
can be bounded by one times an enlarged polynomial, using
\eqref{eq:principal-delta-small}. Finally
\[
 R_9-R_6=J e^{\ell_6}(e^{\ell_9-\ell_6}-1)
\]
and the Leibniz rule prove \eqref{eq:principal-density-jets}.
All operations involve at most $n$ derivatives, so the number
and degree of polynomial factors are fixed, independently of $N$.

\subsection{The good-coordinate region and the Gaussian tail}
On $\|z\|_\infty<v$, the cutoff jets equal those of the
constant one. The difference of the two integrands consequently
has jet norm at most $\tau\delta P_n(N,Q,S)$.
On the rest of the compact ball, the product rule and
\eqref{eq:principal-density-jets} bound its jets by
$\tau P_n(N,Q,S)$, after enlarging the polynomial once more.
Outside that ball the integrands vanish on a parameter
neighborhood, so their jets vanish too.

The reference Gaussian obeys the dimension-explicit tail estimate
\begin{equation}\label{eq:principal-outage}
 \gamma_{C_0}\{\|z\|_\infty\ge v\}
                                  \le2N e^{-v+M/2}.
\end{equation}
For a single coordinate, Markov's inequality applied to $e^{z_i}$
and $e^{-z_i}$ gives
$\mathbb P(|z_i|\ge v)\le2e^{-v+(C_0)_{ii}/2}$.
Sum over the at most $N$ coordinates and use $(C_0)_{ii}\le M$.
No independence of these coordinates is required.

Differentiation under the integral is justified even though $h$
need not be continuous. On a fixed compact coordinate ball the
smooth density-cutoff factors and all their derivatives through
order $n$ are bounded on a smaller parameter neighborhood.
Multiplication by the fixed measurable $h$ supplies an integrable
constant bound. Difference quotients are dominated using the
fundamental theorem of calculus along a short parameter segment.
Apply dominated convergence, successively through order $n$, as
in Section~\ref{sec:gaussian-differentiation}. Integration of the
two pointwise bounds and \eqref{eq:principal-outage} gives
\begin{equation}\label{eq:principal-comparison-final}
 \left\|\int(R_6\chi_6-R_9\chi_9)h\,d\gamma_{C_0}
       \right\|_{J^n_0}
 \le\tau P_n(N,Q,S)\bigl(\delta+2N e^{-v+M/2}\bigr).
\end{equation}
Finite forest and tag averages preserve this estimate, because
their weights are nonnegative and sum to one. The insertion and
the reference transition remain paired before taking these averages.

\subsection{How the comparison is paid}
Write $\beta_f(x)$ for the fine inverse coupling, and use the
upper-exponent convention of Section~\ref{sec:nonlinear-payment-selection}.
Suppose $N,Q,S$ have upper exponent zero, $\delta$ has upper
exponent $-1/2$, and $e^{-v}$ has every negative upper exponent.
Then the multiplier of $\tau$ in
\eqref{eq:principal-comparison-final} has upper exponent $-1/2$.
Indeed a fixed polynomial in exponent-zero quantities still has
upper exponent zero. Its product with $\delta$ has exponent
$-1/2$, whereas its product with $e^{-v}$ can be assigned any
negative exponent. A finite sum preserves the larger exponent.
The smallness requirement \eqref{eq:principal-delta-small}
therefore holds beyond a finite post-length threshold.

Here are the exact scales that produce the comparison data in the
construction. The local support diameter is at most
$D_x=B_x(4B_x+2)+2B_x$, and its read dimension is at most
\[
             N(x)=(2D_x+1)^4L^4\,4\dim\mathfrak g.
\]
This has upper exponent zero for fixed $L$.
The local direction bound is
$2\theta_f(x)(1+D_x)^p$, of upper exponent $-1/2$.
Let $J_0(x)$ bound the unscaled displacement of the two means,
and $V_0(x)$ the covariance difference, including their jets.
The construction requires the bounds
\begin{equation}\label{eq:principal-primitive-comparison}
 J_0:\ -1,\qquad V_0:\ -1/2,
\end{equation}
where the numbers denote upper exponents. Here $\beta$ is the
actual fine inverse coupling of the incoming representation,
within its prescribed window about $\beta_f(x)$.
Gaussian coordinates divide the mean displacement by
$g=\beta^{-1/2}$, so a common
jet difference bound is
\[
             \delta(x)=|V_0(x)|+g^{-1}|J_0(x)|.
\]
Since $g^{-1}$ has upper exponent $1/2$, this has exponent
$-1/2$. Absolute mean and coordinate bounds, covariance inverse
bounds, and the fixed-order Jacobian bounds are exponent zero
after this rescaling. The support radius in Gaussian coordinates
is bounded by
$g^{-1}(5a\theta_f+|c|\theta_f)$, also exponent zero.
The product cutoffs have polynomial fixed-order jet bounds: in a
derivative of order at most $n$, at most $n$ factors are
differentiated; the others have magnitude at most one. There are
at most $N^n$ choices of these factors. The one-factor derivatives
are bounded by the fixed smooth cutoff and its rescaled aperture.

Choose the good-coordinate radius
\[
 v(x)=\frac{\varepsilon}{128(\dim\mathfrak g+1)}\,\mathcal M_f(x),
\]
where $\mathcal M_f$ is the fine monitor. The monitor estimate
gives every negative upper exponent for $e^{-v}$.
The construction's aperture comparison is
$\theta_f/g\ge\mathcal M_f/2$. Eventually
$\delta\le\varepsilon\mathcal M_f/[128(\dim\mathfrak g+1)]$;
hence
\[
 g\dim\mathfrak g\,(v+\delta)<\varepsilon\theta_f/8.
\]
Each reference coordinate then lies strictly inside its cutoff
plateau when $\|z\|_\infty<v$. Continuity gives the required
parameter neighborhood; the cutoff implication in the common
integral gives the other plateau as well. This proves the tail
and numerical purchases directly, once the primitive mean and
covariance comparison \eqref{eq:principal-primitive-comparison}
and their absolute jet bounds have been supplied. Those primitive
estimates concern the actual sector and reference covariances;
the value-only remote mean bound does not by itself prove them.
